\section{Breadth-First Search (BFS) Traversal}
\label{sec:graph-bfs}

\problemstatement{Given a graph (as an adjacency list) and a starting
node, visit every node reachable from the start exactly once,
exploring all nodes at distance $1$ before any node at distance $2$,
and so on.}

\begin{patternbox}
\emph{``Visit level by level, nearest first''} is solved with a
\textbf{queue}: push the start node, then repeatedly pop the front,
visit it, and push every unvisited neighbor onto the back -- the
first-in-first-out order of a queue naturally processes nodes in
increasing distance from the start.
\end{patternbox}

\subsection*{Algorithm}

\begin{enumerate}
  \item Maintain a $\textit{visited}$ array; mark the start node
        visited and push it onto a queue.
  \item While the queue is non-empty: pop the front node, add it to
        the result; for every unvisited neighbor, mark it visited
        \textbf{immediately} (before it is popped) and push it onto
        the queue.
  \item To traverse the whole graph, repeat for every node not yet
        visited.
\end{enumerate}

\begin{ideabox}[Mark Visited on Push, Not on Pop]
A subtle but important detail: a node must be marked visited the
moment it is \textbf{pushed} onto the queue, not when it is popped.
If marking were delayed until pop time, the same node could be pushed
multiple times by different neighbors before any of those pushes is
processed, wasting work and (in graphs with cycles) potentially
causing incorrect results.
\end{ideabox}

\subsection*{C++ implementation}

\begin{lstlisting}
#include <bits/stdc++.h>
using namespace std;

vector<int> bfsFromSource(int src, vector<vector<int>>& adj, vector<bool>& visited) {
    vector<int> result;
    queue<int> q;
    visited[src] = true;
    q.push(src);

    while (!q.empty()) {
        int node = q.front();
        q.pop();
        result.push_back(node);

        for (int neighbor : adj[node]) {
            if (!visited[neighbor]) {
                visited[neighbor] = true;
                q.push(neighbor);
            }
        }
    }
    return result;
}

vector<int> bfsTraversal(int n, vector<vector<int>>& adj) {
    vector<bool> visited(n, false);
    vector<int> result;

    for (int i = 0; i < n; i++) {
        if (!visited[i]) {
            vector<int> part = bfsFromSource(i, adj, visited);
            result.insert(result.end(), part.begin(), part.end());
        }
    }
    return result;
}

int main() {
    int n = 5;
    vector<vector<int>> adj(n);
    auto addEdge = [&](int u, int v) { adj[u].push_back(v); adj[v].push_back(u); };
    addEdge(0, 1); addEdge(0, 2); addEdge(1, 3); addEdge(2, 4);

    for (int x : bfsTraversal(n, adj)) cout << x << " ";
    cout << endl; // 0 1 2 3 4
    return 0;
}
\end{lstlisting}

\subsection*{Dry run}

\dryrun{Take the same graph as Section~\ref{sec:graph-dfs}: edges
$0{-}1$, $0{-}2$, $1{-}3$, $2{-}4$, starting from node $0$.}

Queue: $[0]$ (visited $=\{0\}$). Pop $0$, record; push unvisited
neighbors $1,2$: queue $[1,2]$ (visited $=\{0,1,2\}$). Pop $1$,
record; push unvisited neighbor $3$: queue $[2,3]$. Pop $2$, record;
push unvisited neighbor $4$: queue $[3,4]$. Pop $3$, record; no new
neighbors. Pop $4$, record; no new neighbors. Final order:
$0,1,2,3,4$ -- notice this differs from DFS's $0,1,3,2,4$: BFS visits
both of $0$'s direct neighbors ($1$ and $2$) before descending to
either of their children, exactly the ``nearest first'' property.

\begin{complexitybox}
\textbf{Time Complexity.} $O(V+E)$ -- identical bound to DFS, for the
same reason (every vertex and edge examined once).
\textbf{Space Complexity.} $O(V)$ for the \textit{visited} array and
the queue (which can hold up to all $V$ nodes in the worst case, e.g.\
a star graph).
\end{complexitybox}

\begin{takeawaybox}
DFS and BFS visit the exact same set of reachable nodes and share the
same $O(V+E)$ complexity -- the only difference is traversal
\textbf{order}, governed entirely by swapping a stack for a queue.
Whenever a problem specifically needs shortest paths in an unweighted
graph or strictly level-by-level processing, that need points directly
at BFS; everything else defaults naturally to whichever is more
convenient to implement (usually DFS, via plain recursion).
\end{takeawaybox}
